% sec_concentration.tex
\section{Концентрация и упаковка Амеров в СЭ}

Приведённые выше доводы о величине среднего значения длины свободного пробега в Свободном эфире, или физическом вакууме, \(\lambda_{\text{сэ}} \approx d_a/2\), а также идея Платона о тетраэдральной компоновке применительно к частицам эфира — Амерам — позволяют оценить величину пространства наблюдаемой Вселенной, занимаемого свободным эфиром. \textbf{Карл Фридрих Гаусс} доказал, что самая высокая плотность упаковки для равных сфер, которая может быть достигнута простой регулярной упаковкой, равна \(\pi/(3\sqrt{2}) \approx 0,74048\), а \textbf{Иоганн Кеплер} выдвинул гипотезу о том, что эта упаковка имеет наивысшую плотность среди всех возможных упаковок сфер, регулярных и нерегулярных. Эту гипотезу методом полной индукции в 2017 году доказал американский математик \textbf{Thomas Callister Hales}.

\begin{figure}[h]
\centering
\includegraphics[width=0.55\textwidth]{images/Упаковка Гаусс.png}
\caption{Упаковка Гаусса.}
\label{fig:gauss_packing}
\end{figure}

Изменение длины свободного пробега Амеров в физическом вакууме в пределах \(0 < \lambda_{\text{сэ}} < d_a\) (условие сохранения целостности вихревого эфирного объекта) и хаотичность их движения условно позволяют предполагать взаимное равенство расстояний между всеми Амерами в каждый момент времени, равное среднему значению длины свободного пробега \(d_a/2\), как изображено на рис. \ref{fig:gauss_packing} и \ref{fig:amers_arrangement}. Рассмотрим единичный объём пространства, упакованный равными сферами диаметром \(D = \frac{3}{2}d_a\). Объём сферы \(V = \frac{1}{6}\pi D^3\), выраженный через диаметр Амера \(d_a\), будет иметь вид \(V = \frac{9}{16}\pi d_a^3\). Согласно теореме Иоганна Кеплера наивысшая плотность упаковки сфер в пространстве достигает \(74,05\%\) или \(\pi/(3\sqrt{2})\), что относится в нашем случае к единице объёма пространства, занимаемого Свободным эфиром. Разделив этот объём на объём сферы \(V = \frac{9}{16}\pi d_a^3\), получим величину концентрации Амеров в единице объёма пространства:
\[
n_{\text{сэ}} = \frac{\pi/(3\sqrt{2})}{\frac{9}{16}\pi d_a^3} = \frac{16}{27\sqrt{2}\,d_a^3}.
\]

Определим, какую часть пространства составляет суммарный объём Амеров Свободного эфира:
\[
V_a \cdot n_{\text{сэ}} = \frac{1}{6}\pi d_a^3 \cdot \frac{16}{27\sqrt{2}\,d_a^3} = \frac{8\pi}{81\sqrt{2}} \approx 0.2194 \; (22\%),
\]
где \(V_a\) — объём Амера. Таким образом, Амеры свободного эфира, изотропно заполняя пространство наблюдаемой Вселенной, физически совокупно составляют \textbf{22\%} его объёма, без учёта Амеров, входящих в состав вихревых эфирных объектов весомой материи. Сравним полученный результат с данными, приведёнными В.А. Ацюковским в его «Началах эфиродинамического естествознания» (Москва, 2009, кн.2):
\[
V_a n_{\text{сэ}} = 5,1\cdot10^{-134} \times 5,8\cdot10^{102} = 29,58\cdot10^{-32}.
\]
Нетрудно заметить существенное отличие результатов. Единица объёма пространства — это куб с длиной ребра 1 м. Представим нечто сопоставимое с этими размерами, например, работающий двигатель внутреннего сгорания. Какую роль свободный эфир, составляющий \(29,58\cdot10^{-32}\) от объёма 1 м³, играет в осуществлении и функционировании этой материальной системы? По мнению автора, \textbf{никакую}. Этот двигатель, по сути, представляет пространственную, функционально упорядоченную совокупность объектов весомой материи вихревого происхождения, при этом объём образующих эту совокупность Амеров, как это ни парадоксально, меньше, чем объём составляющих это пространство (1 м³) Амеров свободного эфира, а это \(0,22\) м³ в совокупности. Как показывают цифры, в представляемой модели осуществляется взаимно обусловленное взаимодействие двух фазовых эфирокинетических систем одного пространственного порядка, где, напомню, под фазовыми состояниями подразумеваются различные организационные формы движения частиц эфира-Амеров, а именно хаотическая и направленная (вихревая).

Полученный результат \(0,2194\) (22\%), в случае правильности выдвинутых предположений, имеет интересное следствие, которое относится к величине совокупного объёма частиц в единице объёма пространства \(V_a n_{\text{сэ}}\). Величина этого произведения должна удовлетворять неравенству \(0 < V_a n_{\text{сэ}} \le 1\), что означает физическую невозможность отсутствия Свободного эфира или полное заполнение им всего пространства, а значит, \(V_a\) и \(n_{\text{сэ}}\) имеют величины одного порядка и обратно пропорциональны друг другу, что является следствием их прямой зависимости от значения \(d_a^3\) и «псевдожидкого» фазового состояния свободного эфира.
